\documentclass[aspectratio=169,10pt]{beamer}

% ============================================================
% Packages
% ============================================================
\usepackage[utf8]{inputenc}
\usepackage[T1]{fontenc}
\usepackage{lmodern}
\usepackage{amsmath, amssymb}
\usepackage{booktabs}
\usepackage{graphicx}
\usepackage{ragged2e}
\usepackage{tcolorbox}
% \tcbuselibrary{skins} % requires tikzfill (missing from TeX Live basic)
\usetikzlibrary{arrows.meta}
\usepackage[style=authoryear, backend=biber, maxcitenames=2, maxbibnames=99, uniquename=false, uniquelist=false]{biblatex}
\addbibresource{references.bib}
\usepackage{multicol}
\setbeamertemplate{bibliography item}{}
\ExecuteBibliographyOptions{giveninits=true, isbn=false, url=false, doi=false}
\DeclareNameAlias{default}{family-given}
\DeclareNameAlias{sortname}{family-given}
\renewcommand*{\bibfont}{\scriptsize}
\setlength{\bibitemsep}{2pt}
\renewbibmacro{in:}{}
\setbeamercolor{bibliography entry author}{fg=black}
\setbeamercolor{bibliography entry title}{fg=black}
\setbeamercolor{bibliography entry location}{fg=black}
\setbeamercolor{bibliography entry note}{fg=black}
% References not yet in references.bib are shown in red
\newcommand{\missingref}[1]{\textcolor{red}{#1}}

% ============================================================
% Colours (same as the written proposal)
% ============================================================
\definecolor{boxframe}{RGB}{150,58,78}
\definecolor{boxback}{RGB}{248,238,239}
\definecolor{darkgreen}{RGB}{0,100,0}
\definecolor{textgrey}{RGB}{110,110,110}

% ============================================================
% Beamer look
% ============================================================
\usetheme{default}
\setbeamertemplate{navigation symbols}{}
\setbeamersize{text margin left=0.6cm, text margin right=0.6cm}
\setbeamercolor{frametitle}{fg=boxframe}
\setbeamerfont{frametitle}{series=\bfseries, size=\Large}
\setbeamercolor{structure}{fg=boxframe}
\setbeamertemplate{itemize item}{\color{boxframe}--}
\setbeamertemplate{itemize subitem}{\color{boxframe}\textbullet}
\setbeamertemplate{enumerate item}{\color{boxframe}\insertenumlabel.}

% Header: frame title with a coloured rule underneath
\setbeamertemplate{frametitle}{%
  \vspace{0.4cm}%
  {\usebeamerfont{frametitle}\usebeamercolor[fg]{frametitle}\insertframetitle\par}
  \vspace{0.1cm}
  {\color{boxframe}\rule{\linewidth}{0.6pt}}
}

% Footer: course | title | page
\setbeamertemplate{footline}{%
  \hspace{0.4cm}{\scriptsize\color{textgrey}ENSAE -- IP Paris \;|\; Public Economics \;|\; \insertshorttitle}
  \hfill
  {\scriptsize\color{textgrey}\insertframenumber/\inserttotalframenumber}\hspace{0.4cm}
  \vspace{0.25cm}
}

% ============================================================
% Helpers
% ============================================================
% Grey placeholder text, to replace with content
\newcommand{\placeholder}[1]{{\color{textgrey}\itshape #1}}
% Coloured box, same style as the written proposal
\newtcolorbox{keybox}[1][]{%
  colback=boxback, colframe=boxframe, boxrule=0.6pt, arc=2pt,
  left=4pt, right=4pt, top=2pt, bottom=2pt,
  fonttitle=\bfseries\small, title={#1}
}

\newcommand{\blocklabel}[1]{\vspace{0.15cm}{\color{boxframe}\textbf{#1}}\par\vspace{0.05cm}}

% ============================================================
% Title information
% ============================================================
\title[\placeholder{Who pays when fathers don't?}]{\placeholder{Who pays when fathers don't?}}
\subtitle{\placeholder{Default intermediation of child support in France}}
\author{\placeholder{Eve Samani}, \placeholder{Isaure Pillet}, \placeholder{Antony Albergne}, \placeholder{Lucie Lang}}
\institute{Public Economics -- Enrico Rubolino \\ ENSAE -- IP Paris, 3A}
\date{Fall 2026}

% Title page layout
\setbeamertemplate{title page}{%
  \vfill
  \begin{center}
    {\color{boxframe}\rule{0.7\linewidth}{0.6pt}}\par
    \vspace{0.4cm}
    {\usebeamercolor[fg]{frametitle}\Huge\bfseries\inserttitle\par}
    \vspace{0.25cm}
    {\large\color{textgrey}\insertsubtitle\par}
    \vspace{0.4cm}
    {\color{boxframe}\rule{0.7\linewidth}{0.6pt}}\par
    \vspace{0.6cm}
    {\normalsize\insertauthor\par}
    \vspace{0.3cm}
    {\small\color{textgrey}\insertinstitute\par}
    \vspace{0.2cm}
    {\small\color{textgrey}\insertdate\par}
  \end{center}
  \vfill
}

\begin{document}

% ============================================================
% TITLE SLIDE (no footer, not counted in the slide numbers)
% ============================================================
\begin{frame}[plain, noframenumbering]
  \titlepage
\end{frame}



% ============================================================
% SLIDE 1 -- Motivation and Literature
% ============================================================

% ============================================================
% SLIDE 1 -- Motivation and Literature
% ============================================================

\begin{frame}{Motivation and Literature}

\vspace{-0.35cm}

\begin{columns}[T,onlytextwidth]
% ============================================================
% LEFT COLUMN -- MOTIVATION
% ============================================================
\begin{column}{0.48\textwidth}

  
\blocklabel{Motivation}
\footnotesize

\begin{itemize}
    \setlength{\itemsep}{0.09cm}
    \setlength{\topsep}{0pt}
    \justifying

    \item \textbf{Gendered costs of divorce:}
    Women's living standards fall by
    \textbf{\textcolor{boxframe}{19\%}}
    one year after divorce, compared with
    \textbf{\textcolor{boxframe}{2\%}}
    for men (Bonnet et al., 2021).

    \item \textbf{Unreliable child support:}
    Child support represents
    \textbf{\textcolor{boxframe}{18\%}}
    of single-parent families' resources,
    yet \textbf{\textcolor{boxframe}{30--40\%}}
    of payments were fully or partially unpaid.

    \item \textbf{Policy reform:}
    France progressively introduced
    \textbf{automatic public intermediation}
    of child support (March 2022;
    January 2023) to improve payment reliability.

\end{itemize}

\vspace{0.03cm}
{\tiny\color{textgrey}
Sources: Bonnet et al. (2021); CNAF (2023);
Ministry of Justice (2022).
}
\end{column}

% ============================================================
% RIGHT COLUMN -- RELATED LITERATURE
% ============================================================
\begin{column}{0.48\textwidth}

  \blocklabel{Related literature}
  \footnotesize

  \begin{itemize}
    \setlength{\itemsep}{0.09cm}
    \setlength{\topsep}{0pt}
    \justifying

    \item \textbf{Post-Separation Income \& Support}:
    \textcite{bonnet2021flip} 

    \item \textbf{Agency Failures \& Enforcement}:
    \textcite{weiss1993transfers};
    \textcite{chiappori2007divorce};
    \textcite{freeman2001dunning}.

    \item \textbf{Behavioral Responses \& IFPA}:
    \textcite{rossin2018parental};
    \textcite{cancian2013discouraging};
    \textcite{titlied2024} \textit{(French IFPA)}.

  \end{itemize}

\end{column}

\end{columns}


% ============================================================
% BOTTOM BOXES -- ALIGNED HORIZONTALLY
% ============================================================
\vspace{-0.25cm}

\begin{columns}[T,onlytextwidth]

% LEFT BOX -- RESEARCH QUESTION
\begin{column}{0.48\textwidth}
    \begin{keybox}[Research question]
        \footnotesize
        How does the shift to default public
        child-support intermediation affect
        payment compliance and mothers'
        economic security after separation?
    \end{keybox}
\end{column}

% RIGHT BOX -- OUR CONTRIBUTION
\begin{column}{0.48\textwidth}
    \begin{keybox}[Our contribution]
        \footnotesize
        France's phased reform provides a rare
        opportunity to investigate how public
        enforcement of private transfers shapes
        women's economic security after separation.
    \end{keybox}
\end{column}

\end{columns}
\end{frame}



% SLIDE 2 -- Data and Methodology
% ============================================================
\begin{frame}{Data and Methodology}
\footnotesize
\begin{columns}[T]
  % ---------- Left column: data ----------
  \begin{column}{0.41\textwidth}
    \blocklabel{Data}
    \begin{itemize}\setlength{\itemsep}{1pt}
      \item \textbf{Source}: ALLSTAT-FR6 (CNAF, via CASD), panel 2017--2025
      \item \textbf{Sample}: separated parents, $\geq 1$ child, child support due, $\pm 12$ months around each cutoff
      \item \textbf{Cohort}: month of separation (\texttt{NBMSIFAM})
      \item \textbf{Groups}: married vs.\ unmarried before separation (\texttt{SITFAM})
      \item \textbf{Outcomes $Y$}: child support received, ASF, RSA, prime d'activit\'e, earnings at 6, 12, 24 months
      \item \textbf{Covariates}: age, sex, no.\ of children, pre-separation income, municipality
      \item \textbf{Limitation}: no IFPA status or title date
    \end{itemize}
  \end{column}

  % ---------- Right column: methodology ----------
  \begin{column}{0.57\textwidth}
    \blocklabel{Methodology: difference-in-discontinuities}
    \begin{keybox}[{Specification, at each cutoff $c^*$}]
      \begin{align*}
      Y_i =\; & \underbrace{\textcolor{red}{\beta}\,\text{Post}_i \times \text{Treated}_i}_{\textcolor{red}{\text{effect of default IFPA}}}
      + \underbrace{\tau\,\text{Post}_i}_{\text{common jump}}
      + \underbrace{\theta\,\text{Treated}_i}_{\text{group gap}} \\
      & + \underbrace{f_g(c_i - c^*)}_{\text{cohort trends}}
      + \underbrace{X_i'\gamma}_{\text{controls}}
      + \varepsilon_i
      \end{align*}
      $\textcolor{red}{\beta}$: effect of \emph{default} vs.\ opt-in intermediation (intent-to-treat)
    \end{keybox}
    \begin{itemize}\setlength{\itemsep}{1pt}
      \item \textbf{Assumptions}: (i) absent the reform, both groups jump equally at the cutoff; (ii) no sorting in separation timing or divorce type; (iii) no interaction with concurrent shocks
      \item \textbf{Checks}: no bunching in separations or jump in group shares; covariate balance; placebo cutoffs; staggered DiD
      \item \textbf{Caveat}: married $\neq$ judicial divorce $\Rightarrow$ conservative estimates
    \end{itemize}
  \end{column}
\end{columns}
\end{frame}

% ============================================================
% SLIDE 2bis -- Identification: two cutoffs, two groups
% ============================================================
\begin{frame}{Identification: Two Cutoffs, Two Groups}
\vspace{0.6cm}
\centering
\resizebox{0.95\linewidth}{!}{%
\begin{tikzpicture}[x=1.45cm, y=0.85cm, font=\small]
  \def\cutA{2.33}\def\cutB{4}\def\asf{3.67}\def\xend{6}
  \node[anchor=east, align=right] at (-0.2,2) {\textbf{Judicial divorces}\\[-2pt]{\scriptsize\color{textgrey}(married, divorced by a judge)}};
  \node[anchor=east, align=right] at (-0.2,0.8) {\textbf{Other separations}\\[-2pt]{\scriptsize\color{textgrey}(cohabiting, PACS, notarial divorces)}};
  \fill[textgrey!25] (0,1.75) rectangle (\cutA,2.25);
  \fill[boxframe!75] (\cutA,1.75) rectangle (\xend,2.25);
  \fill[textgrey!25] (0,0.55) rectangle (\cutB,1.05);
  \fill[boxframe!75] (\cutB,0.55) rectangle (\xend,1.05);
  \node[textgrey!80!black] at (\cutA/2,2) {\scriptsize opt-in};
  \node[white] at ({(\cutA+\xend)/2},2) {\scriptsize\bfseries default IFPA};
  \node[textgrey!80!black] at (\cutB/2,0.8) {\scriptsize opt-in};
  \node[white] at ({(\cutB+\xend)/2},0.8) {\scriptsize\bfseries default IFPA};
  \draw[-{Stealth}, thick] (0,0) -- (\xend+0.3,0) node[right] {\scriptsize cohort $c$};
  \foreach \x/\l in {0/2021, 2/2022, 4/2023, 6/2024}
    \draw (\x,0.08) -- (\x,-0.08) node[below] {\scriptsize \l};
  \draw[boxframe, very thick, dashed] (\cutA,-0.1) -- (\cutA,2.6) node[above] {\scriptsize\bfseries Mar.\ 2022};
  \draw[boxframe, very thick, dashed] (\cutB,-0.1) -- (\cutB,2.6) node[above] {\scriptsize\bfseries Jan.\ 2023};
  \draw[textgrey, thin, dotted] (\asf,-0.1) -- (\asf,1.4);
  \node[textgrey, anchor=north] at (\asf,-0.35) {\tiny ASF uprating};
  \node[draw=boxframe, fill=boxback, rounded corners, align=left, anchor=north, text width=5.2cm]
    at (0.9,-0.75) {\scriptsize \textbf{Cutoff 1}: treated = judicial divorces\\ control = other separations};
  \node[draw=boxframe, fill=boxback, rounded corners, align=left, anchor=north, text width=5.2cm]
    at (5.0,-0.75) {\scriptsize \textbf{Cutoff 2}: treated = other separations\\ control = judicial divorces (already treated)};
\end{tikzpicture}}
\end{frame}


% ============================================================
% SLIDE 3 -- Expected Results
% ============================================================

\begin{frame}{Expected Results}

\noindent{\small \color{textgrey}Benchmark: after a divorce or PACS dissolution, women's living standards rely on two cushions, child support and public transfers (Bonnet et al, 2021). The reform shifts the burden between them.}

\vspace{0.2cm}
\begin{itemize}
  \setlength{\itemsep}{0.2cm}
  \item \textbf{H1: Child support received rises}, mainly where fathers \emph{can} pay.

  
  \item \textbf{H2 : Living standards : heterogeneous effects.}
  \begin{itemize}
      \item RSA recipients already receiving support: little change, as it is deducted one-for-one. 
      \item Mothers above benefit thresholds: net income gain.
      \item Mothers whose unpaid support was still counted as income: gain.
  \end{itemize}
  \item \textbf{H3: Labour supply: ambiguous.} Income effect vs\ stabilisation effect (more continuous employment).
\end{itemize}

\vspace{0.15cm}
\begin{keybox}[Who pays when fathers don't?]
  \centering
  Default intermediation $\Rightarrow$ higher compliance $\Rightarrow$ \textbf{better living standards for mothers?}\\[0.05cm]
  {\small Likely for mothers who bore defaults alone, for the poorest, it may only change \emph{who} pays.}
\end{keybox}

\end{frame}

% ============================================================
% REFERENCES (appendix, not counted in the slide numbers)
% ============================================================
\begin{frame}[noframenumbering]{References}
\vspace{-0.3cm}
\scriptsize
\begin{multicols}{2}
  \setlength{\parskip}{3pt}
  \printbibliography[heading=none]

  % Not yet in references.bib (shown in red)
  \everypar{\hangindent=1.2em\hangafter=1}
  \missingref{Cancian, M., Heinrich, C. J., and Chung, Y. (2013). ``Discouraging disadvantaged fathers' employment: an unintended consequence of policies designed to support families''. \textit{Journal of Policy Analysis and Management} 32.4, pp.\ 758--784.}

  \missingref{Freeman, R. B. and Waldfogel, J. (2001). ``Dunning delinquent dads: the effects of child support enforcement policy on child support receipt by never married women''. \textit{Journal of Human Resources} 36.2, pp.\ 207--225.}

  \missingref{Garfinkel, I. and Klawitter, M. M. (1990). ``The effect of routine income withholding of child support on AFDC participation and costs''. \textit{Journal of Policy Analysis and Management} 9.2, pp.\ 155--177.}

  \missingref{HCFEA (2025). \textit{Note sur l'interm\'ediation financi\`ere des pensions alimentaires}. Haut Conseil de la famille, de l'enfance et de l'\^age.}

  {\color{textgrey}Data sources (motivation): CNAF (2023); INSEE (2024); Minist\`ere de la Justice (2022).}
\end{multicols}
\end{frame}

\end{document}